\section{Alternating series and integral tests}
\label{sec:alternating-integral-tests}
In this section, $\log$ means the natural logarithm. A series is
\emph{conditionally convergent} if it converges but its series of
absolute values diverges. Changing or removing finitely many terms
does not change convergence, by Exercise~14.9. We use this fact when
a comparison or monotonicity estimate holds only after a fixed index.

The alternating-series theorem requires nonnegative magnitudes that
decrease to zero. An alternating sign by itself does not guarantee
convergence. An integral test requires a nonnegative decreasing
function matching the terms; the precise statement and proof are
given in Exercise~15.8 below. As in the textbook's Section~15, the
usual elementary integral and derivative formulas are available.

\exercise{15.1}{Two alternating series with very different magnitudes}
\goal Decide whether the series in parts (a) and (b) converge,
and justify each answer.
\idea Check the magnitudes before applying the alternating-series
theorem. In particular, a convergent series must have terms tending
to zero, regardless of their signs.

\step{(a) The series $\sum_{n=1}^{\infty}(-1)^n/n$.}
Put $b_n=1/n$. For every positive integer $n$, we have $b_n>0$.
Also $n+1>n>0$, so taking reciprocals reverses the inequality:
\[
 b_{n+1}=\frac{1}{n+1}<\frac1n=b_n.
\]
Finally, $b_n\to0$. Indeed, for $\varepsilon>0$, choose an integer
$N>1/\varepsilon$. If $n>N$, then
$\abs{b_n-0}=1/n<\varepsilon$.

Thus $\sum_{n=1}^{\infty}(-1)^{n+1}b_n$ converges by the
alternating-series theorem. Multiplication by $-1$ preserves
convergence, by Exercise~14.5(b), and gives the series requested here.
Its series of absolute values is $\sum_{n=1}^{\infty}1/n$, the
divergent harmonic series. Hence the requested series converges
conditionally.

\step{(b) The series $\sum_{n=1}^{\infty}(-1)^n n!/2^n$.}
Let $c_n=n!/2^n>0$. Cancellation of the common factorial and power
gives
\[
 \frac{c_{n+1}}{c_n}
 =\frac{(n+1)!}{2^{n+1}}\frac{2^n}{n!}
 =\frac{n+1}{2}.
\]
For $n\ge3$, this ratio is at least $2$. Repeated multiplication
therefore gives, for $n\ge3$,
\[
 c_n\ge c_3\,2^{n-3}=\frac34\,2^{n-3}.
\]
The right side tends to $+\infty$. In particular $c_n$ does not
tend to zero. If the signed terms $(-1)^nc_n$ tended to zero, their
absolute values $c_n$ would tend to zero too, a contradiction.
The necessary term condition, Corollary~14.5, now proves divergence.
Alternation does not compensate for these growing magnitudes.
\answer{Part (a) converges conditionally. Part (b) diverges because
its terms do not tend to zero.}

\exercise{15.2}{Powers of periodically sampled sine values}
\goal Determine convergence of
\[
 \textup{(a)}\quad\sum_{n=1}^{\infty}
       \left[\sin\left(\frac{n\pi}{6}\right)\right]^n,
 \qquad
 \textup{(b)}\quad\sum_{n=1}^{\infty}
       \left[\sin\left(\frac{n\pi}{7}\right)\right]^n.
\]
The exponent $n$ applies to the whole sine value in each case.
\idea In part (a), the sampled sine values reach $1$ infinitely
often. In part (b), there are only finitely many sampled values,
and every one has absolute value strictly less than $1$.

\step{(a) Exhibit terms that stay equal to $1$.}
For every integer $k\ge0$, choose $n_k=12k+3$. These are strictly
increasing positive indices, and
\[
 \frac{n_k\pi}{6}=2k\pi+\frac{\pi}{2}.
\]
Sine has period $2\pi$ and $\sin(\pi/2)=1$, so
\[
 \left[\sin\left(\frac{n_k\pi}{6}\right)\right]^{n_k}
 =1^{n_k}=1.
\]
The terms of a convergent series must tend to zero, and every
subsequence of a sequence tending to zero must also tend to zero.
This constant subsequence contradicts that requirement. Therefore
the series in part (a) diverges.

\step{(b) Bound all the absolute sine values by one number below $1$.}
The values $\sin(n\pi/7)$ repeat when $n$ is increased by $14$,
because the angle then increases by $2\pi$. Thus all their absolute
values occur in the finite list
\[
 \left\{\abs{\sin\left(\frac{r\pi}{7}\right)}:
                  r=1,2,\ldots,14\right\}.
\]
None of these absolute values equals $1$. Otherwise the angle
$r\pi/7$ would have the form $\pi/2+\ell\pi$ for some integer
$\ell$. After division by $\pi$ and multiplication by $14$, this
would give $2r=7(2\ell+1)$. The left side is even and the right
side is odd, which is impossible.

Let $c$ be the maximum of that finite list. Every entry is less
than $1$, so its attained maximum satisfies $c<1$ as well; also
$c\ge0$. In fact $c>0$, since $\sin(\pi/7)>0$.
For every $n\ge1$,
\[
 \abs{\left[\sin\left(\frac{n\pi}{7}\right)\right]^n}
 =\abs{\sin\left(\frac{n\pi}{7}\right)}^n\le c^n.
\]
The inequality is preserved when taking the positive integer power
because both sides are nonnegative. The geometric series
$\sum_{n=1}^{\infty}c^n$ converges, since $0<c<1$. Comparison proves
that the series of absolute values converges. Hence the series in
part (b) converges absolutely.
\answer{Part (a) diverges by the necessary term condition. Part (b)
converges absolutely by comparison with a geometric series.}

\exercise{15.3}{The logarithmic refinement of the harmonic series}
\goal For a real parameter $p$, prove
\[
 \sum_{n=2}^{\infty}\frac{1}{n(\log n)^p}
 \text{ converges}\quad\Longleftrightarrow\quad p>1.
\]
\idea For positive $p$, use the integral test with the substitution
$u=\log x$. For nonpositive $p$, compare directly with $1/n$.

\step{1. Check positivity and the case $p\le0$.}
For $n\ge2$, $\log n>0$, so every real power $(\log n)^p$ is
defined and positive. Choose an integer $N>e$. For $n\ge N$,
$\log n>1$. If $p\le0$, then $(\log n)^p\le1$, and therefore
\[
 \frac{1}{n(\log n)^p}\ge\frac1n.
\]
The harmonic series diverges. Comparison shows that the series
under consideration diverges to $+\infty$ for $p\le0$.

\step{2. For $p>0$, verify the hypotheses of the integral test.}
Define
\[
 f_p(x)=\frac{1}{x(\log x)^p}\qquad(x\ge2).
\]
The function is positive and continuous on this interval.
Differentiating the product $x^{-1}(\log x)^{-p}$ gives
\[
 \begin{aligned}
 f_p'(x)
 &=-x^{-2}(\log x)^{-p}
   -p x^{-2}(\log x)^{-p-1}\\
 &=-\frac{\log x+p}{x^2(\log x)^{p+1}}<0.
 \end{aligned}
\]
The denominator is positive, and the numerator $\log x+p$ is
positive because $x\ge2$ and $p>0$. Thus the function decreases
and matches the series terms at every integer $n\ge2$.

\step{3. Evaluate the integrals and distinguish the parameter cases.}
For $R\ge2$, put $u=\log x$, so $du=dx/x$. Then
\[
 \int_2^R\frac{dx}{x(\log x)^p}
 =\int_{\log2}^{\log R}u^{-p}\,du.
\]
If $p\ne1$, the power rule gives
\[
 \int_2^R\frac{dx}{x(\log x)^p}
 =\frac{(\log R)^{1-p}-(\log2)^{1-p}}{1-p}.
\]
For $p>1$, one has $1-p<0$, so $(\log R)^{1-p}\to0$ as
$R\to\infty$. The integral tends to the finite positive number
$(\log2)^{1-p}/(p-1)$. The integral test proves convergence.

For $0<p<1$, one has $1-p>0$, so $(\log R)^{1-p}\to+\infty$.
The denominator $1-p$ is positive, and the integral tends to
$+\infty$. The integral test proves divergence.

For $p=1$, integration instead gives
\[
 \int_2^R\frac{dx}{x\log x}
 =\log(\log R)-\log(\log2)\longrightarrow+\infty.
\]
The arguments of both outer logarithms are positive, even though
$\log(\log2)$ itself is negative. This case also diverges.
\answer{Exactly the parameters $p>1$ give convergence. For every
$p\le1$, the nonnegative series diverges to $+\infty$.}

\exercise{15.4}{Four series involving logarithms}
\goal Determine convergence of each of the four requested series.
\idea Compare the first two series with the harmonic series.
For the other two, check the integral-test hypotheses and compute
the relevant integrals.

\step{(a) The series $\sum_{n=2}^{\infty}1/(\sqrt n\log n)$.}
We first justify the eventual inequality $\log x\le\sqrt x$.
Let $g(x)=\sqrt x-\log x$ for $x>0$. Its derivative is
\[
 g'(x)=\frac{1}{2\sqrt x}-\frac1x
       =\frac{\sqrt x-2}{2x}\ge0\qquad(x\ge4).
\]
Also $\log2=\int_1^2dt/t<1$, since $1/t<1$ on $(1,2]$.
Thus $g(4)=2-\log4=2-2\log2>0$. As $g$ is nondecreasing
on $[4,\infty)$, we have $\log x<\sqrt x$ there.

For integers $n\ge4$, both logarithms and square roots are positive,
and the inequality just proved gives
\[
 \sqrt n\log n\le n,
 \qquad
 \frac{1}{\sqrt n\log n}\ge\frac1n.
\]
Taking reciprocals reverses the positive-denominator inequality.
Comparison with the divergent harmonic series proves divergence
to $+\infty$.

\step{(b) The series $\sum_{n=2}^{\infty}(\log n)/n$.}
Choose an integer $N>e$. If $n\ge N$, monotonicity of the
logarithm gives $\log n>\log e=1$. Consequently
\[
 \frac{\log n}{n}\ge\frac1n.
\]
All terms are positive for $n\ge2$. Comparison with the harmonic
series, with the first finitely many indices ignored, proves
divergence to $+\infty$.

\step{(c) The series $\sum_{n=4}^{\infty}1/[n(\log n)(\log\log n)]$.}
For $x\ge4$, we have $\log x>1$. One direct check at the left
endpoint is
\[
 \log2=\int_1^2\frac{dt}{t}>\frac12,
 \qquad \log4=2\log2>1.
\]
Therefore $\log\log x>0$ on this whole interval.
The three factors $x$, $\log x$, and $\log\log x$ are positive,
continuous, and increasing. Their product is increasing, so its
positive reciprocal
\[
 f(x)=\frac{1}{x(\log x)(\log\log x)}
\]
is decreasing and continuous on $[4,\infty)$.

Two substitutions, first $u=\log x$ and then $v=\log u$, give
\[
 \begin{aligned}
 \int_4^R\frac{dx}{x(\log x)(\log\log x)}
 &=\int_{\log4}^{\log R}\frac{du}{u\log u}\\
 &=\int_{\log\log4}^{\log\log R}\frac{dv}{v}\\
 &=\log(\log(\log R))-\log(\log(\log4)).
 \end{aligned}
\]
Every argument of a logarithm in these formulas is positive for
$R\ge4$: in particular $\log\log4>0$. As $R\to\infty$, the
first triple logarithm tends to $+\infty$. The integral test proves
that this series diverges to $+\infty$.

\step{(d) The series $\sum_{n=2}^{\infty}(\log n)/n^2$.}
Let $h(x)=(\log x)/x^2$ for $x\ge2$. This function is positive
and continuous. Differentiation gives
\[
 h'(x)=\frac{1-2\log x}{x^3}<0\qquad(x\ge2),
\]
because $\log x\ge\log2>1/2$, as established in part (c).
Thus the integral test applies. For integration by parts, take
$u=\log x$ and $dv=x^{-2}\,dx$. Then $du=dx/x$ and $v=-1/x$,
so
\[
 \begin{aligned}
 \int_2^R\frac{\log x}{x^2}\,dx
 &=\left[-\frac{\log x}{x}\right]_2^R
      +\int_2^R\frac{dx}{x^2}\\
 &=\frac{\log2+1}{2}-\frac{\log R+1}{R}.
 \end{aligned}
\]
Part (a) gives $0\le(\log R)/R\le1/\sqrt R$ for $R\ge4$.
The right side tends to zero, as does $1/R$. Hence the integral
tends to $(\log2+1)/2$, a finite number. The integral test proves
convergence of the series.
\answer{Parts (a), (b), and (c) diverge to $+\infty$.
Part (d) converges.}

\exercise{15.5}{Why a direct comparison does not establish all positive powers}
\goal Explain why the text used an integral argument to prove the
convergence of $\sum 1/n^p$ for every $p>1$.
\idea A comparison proof needs an independently known convergent
upper series. It cannot assume a nearby case of the very theorem
it is trying to establish.

\step{1. Check which direction of comparison is useful.}
Knowing that $\sum1/n^2$ converges immediately handles $p\ge2$,
since $0\le1/n^p\le1/n^2$ for $n\ge1$. For $1<p<2$, however,
the inequality reverses:
\[
 \frac1{n^p}\ge\frac1{n^2}.
\]
Being larger than the terms of a convergent nonnegative series
does not establish either convergence or divergence. Comparison
with $1/n$ also fails to settle the question: for $p>1$ the terms
$1/n^p$ are smaller than the terms of that divergent series.

\step{2. Explain the missing starting exponent.}
If convergence were already known at some $p_0>1$, comparison
would prove it for every $p\ge p_0$. It would leave the interval
$1<p<p_0$ unresolved. There is no smallest real exponent greater
than $1$: between $1$ and $p_0$ lies $(1+p_0)/2$.
Thus no single known case $p_0>1$ settles the whole range by this
direct comparison. Choosing an unproved exponent $q$ with
$1<q<p$ merely replaces the original convergence problem with
another case of the same problem. The integral calculation supplies
an independent argument for every $p>1$.

\step{3. Comparison can still work after grouping the terms.}
The explanation above does not mean a comparison proof is impossible.
For example, fix $p>1$ and group the indices into blocks
\[
 2^j\le n<2^{j+1}\qquad(j=0,1,2,\ldots).
\]
There are exactly $2^j$ integers in this block, and each satisfies
$n^{-p}\le2^{-jp}$. Hence the sum of the block is at most
\[
 2^j2^{-jp}=2^{j(1-p)}.
\]
The bounding block sums form a geometric series with ratio
$2^{1-p}<1$. Every original partial sum is bounded above by the
sum of this geometric series, since its terms are nonnegative and
its indices occupy finitely many blocks. The original partial sums
are increasing and bounded above, so they converge by Theorem~10.2.
This is a valid alternative proof using an extra grouping argument.
\answer{The integral proof avoids relying on an unproved nearby
$p$-series. A plain comparison with one known exponent does not cover
all $p>1$; a geometric comparison of grouped blocks does.}

\exercise{15.6}{What squaring the terms can change}
\goal Supply the examples in parts (a) and (c), and justify the
nonnegative-term assertion in part (b).

\step{(a) A divergent series whose squared-term series converges.}
Take $a_n=1/n$. Then $\sum a_n$ is the divergent harmonic series,
whereas
\[
 \sum_{n=1}^{\infty}a_n^2
 =\sum_{n=1}^{\infty}\frac1{n^2}
\]
converges by the $p$-series theorem, with $p=2>1$.
These two facts verify both required properties of the example.

\step{(b) Squaring preserves convergence when the original terms are nonnegative.}
Suppose $a_n\ge0$ and $\sum a_n$ converges. Corollary~14.5 gives
$a_n\to0$. There is therefore an integer $N$ with $0\le a_n<1$
whenever $n>N$. Multiplication of $a_n<1$ by the nonnegative
number $a_n$ gives
\[
 0\le a_n^2\le a_n\qquad(n>N).
\]
The comparison test proves convergence of $\sum_{n>N}a_n^2$.
Adding its finitely many omitted real terms preserves convergence,
so the whole squared-term series converges. This is also
Exercise~14.7 with the power $p=2$.

\step{(c) A convergent signed series whose squared-term series diverges.}
Take
\[
 a_n=\frac{(-1)^{n+1}}{\sqrt n}.
\]
Its magnitudes $b_n=1/\sqrt n$ are positive and decrease, because
the positive square root increases with its argument and taking
reciprocals reverses order. They tend to zero: given
$\varepsilon>0$, choose $N>1/\varepsilon^2$; then $n>N$
implies $1/\sqrt n<\varepsilon$.
The alternating-series theorem proves that $\sum a_n$ converges.
However,
\[
 a_n^2=\frac{[(-1)^{n+1}]^2}{(\sqrt n)^2}=\frac1n,
\]
so $\sum a_n^2$ is the divergent harmonic series.
\answer{Use $a_n=1/n$ for part (a) and
$a_n=(-1)^{n+1}/\sqrt n$ for part (c). Part (b) follows from the
eventual bound $0\le a_n^2\le a_n$.}

\exercise{15.7}{A necessary decay rate for decreasing summable terms}
\goal Prove that if $(a_n)$ is a decreasing real sequence and
$\sum a_n$ converges, then $na_n\to0$. Then apply this conclusion
to the harmonic series.
\idea First deduce that every $a_n$ is nonnegative. A late block
of roughly $n/2$ terms is small by the Cauchy criterion, yet every
term in that block is at least $a_n$.

\step{(a), first: derive the sign condition from the actual hypotheses.}
Since $\sum a_n$ converges, Corollary~14.5 gives $a_n\to0$.
The word decreasing means $a_{n+1}\le a_n$; strict decrease is
not required. If $a_j<0$ at any index $j$, then $a_n\le a_j<0$
for all $n\ge j$. Such a tail cannot tend to zero: with
$\varepsilon=-a_j/2>0$, its terms satisfy
$\abs{a_n}\ge-a_j>\varepsilon$. This contradicts $a_n\to0$.
Thus $a_n\ge0$ for every index $n$.

\step{(a), second: use a tail block with an explicit number of terms.}
Fix $\varepsilon>0$. By the Cauchy criterion for the convergent
series, choose an integer $N\ge1$ such that
\[
 q\ge m>N\quad\Longrightarrow\quad
 \abs{\sum_{j=m}^{q}a_j}<\frac{\varepsilon}{2}.
\]
For any integer $n>2N$, set $k=\lfloor n/2\rfloor$.
Here $\lfloor x\rfloor$ is the greatest integer at most $x$.
We have $k\ge N$, so the block starting at $k+1$ begins after $N$.
It contains exactly $n-k$ terms, and $n-k\ge n/2$.
For $k+1\le j\le n$, decreasingness gives $a_j\ge a_n$.
Consequently
\[
 0\le\frac n2 a_n
 \le(n-k)a_n
 \le\sum_{j=k+1}^{n}a_j
 =\abs{\sum_{j=k+1}^{n}a_j}
 <\frac{\varepsilon}{2}.
\]
The equality with the absolute value uses nonnegativity of the
terms. Multiplying the resulting strict upper bound by $2$ gives
$0\le na_n<\varepsilon$ whenever $n>2N$.
This verifies the definition of $na_n\to0$.

\step{(b) Apply the necessary condition to $a_n=1/n$.}
The sequence $1/n$ is decreasing. If the harmonic series converged,
part (a) would imply
\[
 n\left(\frac1n\right)\longrightarrow0.
\]
The expression on the left equals $1$ for every $n$, so its limit
is $1$, a contradiction. Therefore the harmonic series diverges.
Since its terms are nonnegative, its partial sums are increasing;
having no finite limit, they tend to $+\infty$ by the monotone
sequence theorem for unbounded increasing sequences.
\answer{Every decreasing real sequence with a convergent series
satisfies $na_n\to0$. The constant value $n(1/n)=1$ therefore
provides another proof of divergence of the harmonic series.}

\exercise{15.8}{A general integral test and its error bounds}
\goal State and prove a precise integral test of the type described
in Discussion~15.2.
\idea Compare each unit-width rectangle with the integral below
and above it. Nonnegativity makes both the partial sums and the
integrals increasing, so a finite upper bound implies convergence.

\step{1. State the assumptions and the conclusion.}
Let $m\ge1$ be an integer, and let
$f:[m,\infty)\to[0,\infty)$ be continuous and nonincreasing.
Set $a_n=f(n)$ for every integer $n\ge m$.
Define the improper integral by
\[
 \int_m^{\infty}f(x)\,dx
 =\lim_{R\to\infty}\int_m^R f(x)\,dx,
\]
where the value $+\infty$ is allowed. Then
\[
 \sum_{n=m}^{\infty}a_n\text{ converges}
 \quad\Longleftrightarrow\quad
 \int_m^{\infty}f(x)\,dx<+\infty.
\]
If either side is infinite, both the integral and the series diverge
to $+\infty$. If they converge and
$R_N=\sum_{n=N+1}^{\infty}a_n$ for $N\ge m$, then additionally
\[
 \int_{N+1}^{\infty}f(x)\,dx
 \le R_N\le\int_N^{\infty}f(x)\,dx.
\]
Continuity is a sufficient condition ensuring that all the finite
integrals used here exist. The proof also works whenever those
integrals exist, with the same nonnegativity and monotonicity
assumptions.

\step{2. Establish the unit-interval inequalities.}
Fix an integer $k\ge m$. If $k\le x\le k+1$, nonincreasingness
gives $f(k+1)\le f(x)\le f(k)$. Integrating these constant bounds
over the interval of length $1$ yields
\[
 f(k+1)\le\int_k^{k+1}f(x)\,dx\le f(k).
\]
In particular, summing the right-hand inequalities for
$k=m,m+1,\ldots,N$ gives
\[
 \int_m^{N+1}f(x)\,dx\le\sum_{k=m}^{N}f(k).
\]
For the opposite comparison, use
$f(k)\le\int_{k-1}^{k}f(x)\,dx$ for $k\ge m+1$.
Keep the first term $f(m)$ separate and sum to obtain
\[
 \sum_{k=m}^{N}f(k)\le f(m)+\int_m^{N}f(x)\,dx.
\]
At $N=m$, the second sum has no remaining terms and the integral
is zero, so this estimate still holds. Combining the two bounds,
with $S_N=\sum_{k=m}^{N}f(k)$, gives
\[
 \boxed{\int_m^{N+1}f(x)\,dx\le S_N
        \le f(m)+\int_m^{N}f(x)\,dx.}
\]

\step{3. A finite integral bounds the partial sums.}
Suppose the improper integral has finite value $I$.
Since $f\ge0$, one has $\int_m^N f\le I$. Therefore
$0\le S_N\le f(m)+I$ for every $N\ge m$.
Also $S_{N+1}-S_N=f(N+1)\ge0$, so $(S_N)$ is increasing.
The monotone sequence theorem gives a finite limit of $(S_N)$,
which is precisely convergence of the series.

\step{4. A convergent series bounds all the finite integrals.}
Conversely, suppose $S_N\to S<+\infty$. The nonnegative terms
imply $S_N\le S$. If $R\ge m$ is any real number, choose an
integer $N\ge m$ with $R\le N+1$. Nonnegativity and the first
rectangle estimate give
\[
 0\le\int_m^R f(x)\,dx
 \le\int_m^{N+1}f(x)\,dx\le S_N\le S.
\]
Thus all the integrals are bounded above by $S$, not only those
with integer endpoints. They increase with $R$, again because
$f\ge0$. Their supremum is finite, and their limit as
$R\to\infty$ equals that supremum: given a positive tolerance,
choose an endpoint whose integral is within that tolerance of the
supremum; every larger endpoint gives an integral between that
value and the supremum. Hence the improper integral converges.

\step{5. Explain divergence and derive the tail estimates.}
If the integral tends to $+\infty$, the lower bound
$S_N\ge\int_m^{N+1}f$ forces $S_N\to+\infty$.
If instead $S_N\to+\infty$, the upper rectangle bound rearranges
to $\int_m^N f\ge S_N-f(m)$, which forces the integral to
$+\infty$ as well. Nonnegative increasing sums and integrals have
no other type of divergence.

Finally, in the convergent case, apply the same unit-interval
comparisons to the terms with indices $N+1,\ldots,M$, where
$M>N\ge m$:
\[
 \int_{N+1}^{M+1}f(x)\,dx
 \le\sum_{k=N+1}^{M}f(k)
 \le\int_N^M f(x)\,dx.
\]
Letting $M\to\infty$ gives the stated remainder bounds. These
limits are finite since the whole integral and series converge.

The same test applies if the assumptions hold only after some
integer index: start at that index and use Exercise~14.9 to restore
the finite initial segment.
\answer{For a continuous nonnegative nonincreasing function matching
the terms, the series converges exactly when its improper integral
is finite. The two neighboring tail integrals bound the omitted
tail of the series.}
